% T3_rmse.tex - absolute RMSE, model and all three references, array11.
% Sources (verified against both directly, values agree to rounding):
%   Model RMSE: results/table4_protocol_lagged.csv, config_id=C1,
%   array11 (daylight hours, lagged regime, XGBoost, skill vs convex
%   reference - the protocol used throughout this paper; skill column
%   there matches xgb_skill_vs_convex in reference_comparison.csv
%   exactly at all three horizons).
%   Reference RMSEs: results/reference_comparison.csv, array11,
%   rmse_persistence / rmse_climatology / rmse_convex.
% Needs \usepackage{booktabs}.
\begin{table}[t]
  \centering
  \small
  \caption{Absolute RMSE (kW), array 11, lagged regime, validation year
  2014, daylight hours only. Skill is a ratio; this table lets it be
  reconstructed and shows whether each reference is weak or strong in
  absolute terms. Climatology is constant across horizons because it
  uses no recent observation. Persistence is competitive with the
  model at $h=1$ but degrades sharply by $h=6$; climatology and the
  convex reference do not.}
  \label{tab:T3}
  \begin{tabular}{c rrrr}
    \toprule
    $h$ & Model & Persistence & Climatology & Convex \\
    \midrule
    1 & 0.318 & 0.428 & 0.669 & 0.400 \\
    3 & 0.448 & 0.955 & 0.669 & 0.621 \\
    6 & 0.535 & 1.548 & 0.669 & 0.666 \\
    \bottomrule
  \end{tabular}
\end{table}
